\section{Institutional design, implementation and empirical priorities}
\subsection{An ownership contract, not a substitute for the state}
An implementable usownership contract must specify the entity whose equity is allocated, eligible persons, the issue schedule, the unit of contribution, vesting, cash distributions, voting rights and treatment of corporate actions. Shares should attach to the consolidated source of surplus rather than a shell subsidiary from which royalties and valuable assets can be removed. Mergers, buybacks, private recapitalisations and offshore restructurings need explicit rules. Those issues are shared with in-kind equity taxation, not eliminated by a token register.

A consumer, safety evaluator and plugin creator can have different legitimate claims. The score must nevertheless distinguish customer use from paid labour, pre-existing intellectual property and personal-data consent. The baseline's concavity in expenditure limits the advantage of wealthy consumers but does not establish an optimal distribution. Raw attention time should not be an uncapped reward base. Productive evaluations should be audited against outcomes where possible, and appealable decisions should not depend solely on an opaque AI classifier. Sybil resistance, proportional audits and exclusion safeguards are implementation costs to measure, not decorative technical features.

The minimum architecture could be a regulated digital share register with conventional payment accounts and auditable aggregate issuance; a stronger Internet-of-Value architecture would place tokenised claims and programmable settlement on interoperable rails. Distributed ledgers may assist coordination across independent institutions; privacy-preserving proofs may verify an allocation computation without publishing raw behavioural histories. Neither verifies that the underlying contribution data are true. Full public transaction histories are not required and would create disproportionate privacy risks. AI is appropriate for decision support and error detection, not unreviewable monetary or property-right decisions.

\subsection{Money, public services and the European boundary}
The sovereign-money assumption is economically consequential. Moving commercial deposits into sovereign accounts requires a bank-funding and legacy-debt transition; it does not give the government a free claim equal to all existing broad money. Foreign-currency substitution and exchange-rate pressure may make money demand more elastic than in this closed model. A legal ban on other payment instruments does not eliminate alternative stores of value.

For Europe, economic feasibility in a hypothetical sovereign system must be distinguished from present law. Article 123 TFEU restricts central-bank financing of public authorities \citep{tfeu123}. The digital-euro work is not an existing mandate for the programmable demurrage and fiscal financing assumed here \citep{ecbpilot2026,bundesbank2026}. The manuscript therefore proposes neither unilateral national euro issuance nor an assertion that the full architecture can be implemented under current treaties. A legal evaluation of mandatory share issuance, corporate governance, securities, data protection and cross-border capital movements is required in parallel with economic validation.

Public provision must remain explicit. A 7\% public-services case and the 22\% services block in the central model represent different commitments; neither equals the total observed general-government budget, which includes transfers and other items. Moving healthcare or pensions to private expenditure is not an administrative saving. No large bureaucracy saving, subsidy multiplier or automatic tax-compliance advantage is imported from the original exploratory calculations.

\subsection{A staged research and implementation path}
The participant-equity contract can be tested before monetary-regime replacement. A first empirical design would compare actual equity awards with present-value-matched cash rewards, plus a cash reward that accumulates into a financial asset. These arms distinguish incentives from simple payment size and distinguish user equity from asset ownership more generally. Matched eligibility, contribution tasks and transparent consent should be held fixed. Relevant outcomes include productive feedback quality, retention, additional sales, owner returns, participation exclusion, administration cost and portfolio concentration. Attrition and non-participation must be reported, not discarded.

A second stage would estimate how the award schedule affects entry, finance and investment. Both dilution and a matched levy should face common demand shocks and an explicit commitment environment. A third stage would link those estimates to household survey data and a bank/debt/external-sector model. Only then could the monetary bridge be evaluated against debt finance and conventional fiscal adjustment on a comparable macroeconomic basis. This sequencing tests the most distinctive institutional component without requiring a sovereign currency reform to learn whether participants value equity.

A much-later diversification path can preserve the initial design: firm-specific claims, vested portability, voluntary exchange into diversified assets and only subsequently possible sectoral or public pools. A century is an illustrative scenario horizon, not a prescribed institutional timetable. Even within an extended firm-specific phase, households need access to risk management; preserving the participation principle does not require exposing them permanently to one firm's bankruptcy.

\subsection{Novelty, feasibility and the 2027--2037 opportunity window}
The novelty is best stated as an integration across an identifiable research lineage rather than as a claim that every ingredient is new. De la Rosa Esteva's 2022 IoV chapter anticipated a ``Value Deluge'' associated with the large-scale migration of value onto ledgers and linked interoperable IoV infrastructure to the democratisation of finance and property \citep{delarosa2022iov}. ERC-4950, created the same year, explicitly introduced ``Usowners'' as users who can become partial NFT owners and proposed entangled tokens as a linked-control primitive \citep{erc4950}. The present paper adds the missing generalisation: usownership is a recurring firm-level institution with heterogeneous contribution weights, exact dilution accounting, automation-linked accumulation, coverage constraints, fiscal replication tests and a bounded monetary bridge. The contribution is therefore not the first use of partial user ownership; it is the structural-transition model that connects an ownership primitive to the changing distribution of production income.

Feasibility has three distinct layers. Technically, the enabling environment is materially closer than it was in 2022. The ECB reports that European issuers had placed close to EUR~4 billion of DLT-based fixed-income instruments since 2021 and that Eurosystem exploratory work processed roughly EUR~1.6 billion across nine jurisdictions \citep{cipollone2026}. BIS work in 2026 similarly treats tokenisation and trusted programmable money as a plausible next-generation financial architecture rather than as a purely crypto-native experiment \citep{bis2026}. Institutionally, however, tokenised shares remain shares: corporate, securities, competition, privacy and insolvency law continue to govern the claim. Economically, the present simulations show that neither dilution nor monetary finance is free. A pilot can therefore be technologically feasible well before a mandatory national regime is economically or legally justified.

The next decade is relevant because the ownership mechanism is intentionally gradual. At a constant 1\% annual new-share issue, Proposition~\ref{prop:accumulation} implies that users collectively reach only about 9.5\% ownership after ten years; at 2\%, the corresponding share is about 18.0\%. A mechanism intended to cushion a future ownership gap cannot be designed only after that gap has become acute. At the same time, present evidence does not justify treating radical automation by 2037 as certain. \citet{katz2026} estimate modest realised aggregate GDP effects from generative AI through 2025, whereas \citet{fannguyen2026} value observed AI-enabled time savings at a labour-cost equivalent of USD~2.7 trillion, or 3.4\% of GDP, and transformative-AI models permit much larger factor-distribution changes under stronger assumptions \citep{trammell2026}. The 2027--2037 period is therefore best understood as an \emph{institution-design window under deep uncertainty}: production technology and financial rails are both changing, while ownership institutions can be tested before path dependence and concentration become harder to reverse.

\begin{table}[htbp]\centering\small
\caption{Illustrative evidence-gated path for the next decade; dates are design horizons, not forecasts}
\begin{tabularx}{\linewidth}{@{}p{.16\linewidth}p{.27\linewidth}p{.31\linewidth}X@{}}
\toprule
Illustrative horizon & Usownership institution & Programmable-value infrastructure & Evidence gate before scaling \\
\midrule
2027--2030 & Voluntary or contract-based firm pilots; matched equity/cash experiments; vesting and retained-core tests & Regulated digital share registers or tokenised securities; auditable contribution records & Participation coverage, investor response, administration cost, score manipulation and legal treatment \\
2031--2033 & Interoperable firm-specific awards; portability and diversification options; stronger anti-reconcentration rules & Common token/registry standards, automated corporate actions and conditional settlement & Demonstrated net operating-cost advantage and acceptable privacy, oracle, cyber and market-liquidity risks \\
2034--2037 & Conditional wider adoption in sectors where automation materially shifts factor income & Integration with trusted tokenised money or central-bank settlement where legally available & Measured labour-share/displacement effects, money demand, inflation headroom and evidence that ownership income reduces transfer needs \\
\bottomrule
\end{tabularx}

\end{table}

The paper deliberately keeps two time horizons distinct. The first is the 2027--2037 institutional window, in which firm-level pilots, contribution measurement, legal design and programmable-value rails can be tested before any economy-wide adoption. The second is a 50-year structural horizon, illustratively 2027--2077, used to study the cumulative ownership transition rather than to date a forecast. The distinction matters because compounding is slow: absent sales and other issuance, a constant 1\% annual issue produces about 9.5\% user ownership after ten years but about 39.2\% after fifty years. Model year 50 therefore remains the principal long-run comparison for whether usownership can become a material source of household capital income and reduce reliance on MVI and monetary finance; the 100-year extension is retained only as a more remote sensitivity exercise.

\section{Limitations and conclusion}
The main unestimated components are the productivity and investment responses to ownership, the causal value and cost of contribution measurement, relative institutional durability, corporate decision rights, and monetary demand after regime change. The model omits bank credit, sovereign debt, trade, market equity prices, endogenous corporate entry and generational renewal. Its declining wage share and resource ceiling are assumed. These omissions are not minor uncertainty intervals around a forecast; they define the domain of the experiment.

The revised comparison supports a specific novelty claim. Earlier work in this research trajectory had already anticipated the migration of value onto ledgers and the need to preserve ownership across an Internet of Value \citep{delarosa2022iov}; ERC-4950 had already named ``Usowners'' and demonstrated a digital partial-ownership primitive \citep{erc4950}. This paper does not relabel those precursors as new. It generalises them into a macroeconomic institution and subjects that institution to tests absent from the precursors: heterogeneous contribution weights, exact corporate dilution, investor cost, exclusion and reconcentration, matched dividend-levy replication, stock-flow-consistent monetary accounting and transition scenarios. The strongest contribution is thus a bridge between the economics of structural change and a programmable property-rights architecture: an attempt to make redistribution increasingly \emph{pre-distributive}, by changing who accumulates productive claims before profits have to be repeatedly taxed and transferred.

Feasibility is correspondingly conditional but increasingly concrete. The simulations show that gradual participation equity can diffuse legal ownership and reduce future transfer requirements under the stated allocation and payout rules; most of the central hybrid's transfer reduction arises from ownership diffusion rather than the holding-levy monetary layer. A matched dividend levy can reproduce the same cash income under replication assumptions, so the case for usownership must rest on what the legal asset changes: durable claims, cumulative wealth, governance, participation incentives, portability and the possibility that future distributions occur without a fresh fiscal appropriation. Programmable value infrastructure can make recurring fractional awards, vesting and settlement cheaper to execute at scale, and European tokenised-market experiments demonstrate that relevant rails are no longer hypothetical \citep{cipollone2026,eciov2026,bis2026}. Yet technology does not remove dilution, securities law, contributor-verification costs or portfolio risk. The most credible near-term implementation is therefore staged and evidence-gated rather than an immediate economy-wide mandate.

The strategic opportunity lies in timing. If AI and robotics profoundly alter production during 2027--2037, waiting for labour income to collapse before building alternative ownership channels would leave policy dependent on ex post taxation and transfers precisely when the labour tax base may be weakening. If the transition is slower, early pilots still generate evidence at limited macroeconomic cost. Because usownership accumulates gradually, the institution has option value when started early: it can remain small if displacement is modest and scale only if automation materially shifts factor income. The same logic confines seigniorage to a supporting role. Bounded monetary finance can bridge part of the interval before distributed capital income matures, but money demand, inflation, real supply and public commitments remain binding constraints. The paper therefore proposes neither a taxless theorem nor a prediction of a post-work economy. It identifies a feasible research programme for 2027--2037: test whether an Internet-of-Value infrastructure can convert measured participation into durable productive ownership quickly enough to reduce the social cost of a potentially radical AI production transition, while preserving innovation, universal protection and monetary stability. The accompanying 50-year horizon, illustratively extending to 2077, asks the different structural question of whether those initially small ownership flows can compound into a sufficiently broad capital-income base to reduce MVI and seigniorage dependence over a generation-scale transition.

\section*{Data, code and generative-AI disclosure}
All simulated observations are synthetic and reproducible from the accompanying source, configurations and fixed seeds. Official aggregate anchors are separately labelled. The code, tables and figures accompany the manuscript; a public archive identifier has not yet been registered. Generative-AI assistance through ChatGPT (OpenAI) was used for literature discovery, formalisation, code development, computational checks and drafting. Reported numerical results were obtained by executing the supplied Python code; plots were generated from those results, not by synthetic-image generation. Authorship, accountability and approval belong to the submitting human authors. Funding, competing interests and author contributions are provided separately after author confirmation.
